\section{Sobolev 空间}\label{chap:sobolev-spaces}

偏微分方程的分析自然涉及不仅由函数本身的性质, 而且由其导数的性质定义的函数空间. 例如, 在 Problem~\ref{ex:bounded-continuous-derivatives-banach} 中, 引入了 Banach 空间 $C_b^k(\overline\Omega)$. 遗憾的是, 这些空间对于分析 PDE 很不合适. 例如, 若 $\Delta u=f$ 且 $f$ 连续, 一般并不能断言 $u$ 为 $C^2$ 函数. Sobolev 空间要有用得多. 它们的定义方式与 $C_b^k$ 类似, 但由 $L^p$ 取代连续有界函数的角色. 在本节中, $\Omega$ 始终为 $\mathbb{R}^m$ 中的局部 Jordan 可测开集.

本章旨在向读者简要概述 Sobolev 空间. 为便于读者在深入证明之前掌握主要思想, 将材料安排如下.

\begin{itemize}
\item 首先介绍 ``正阶'' Sobolev 空间的基本定义, 属于 (及不属于) 这些空间的典型函数的例子, 以及这些空间的基本性质.
\item 然后用一节讨论 Sobolev 空间的各种重要性质. 将考察以下问题.
\begin{enumerate}
\item 若 $\Omega$ 为整个 $\mathbb{R}^m$, 可以定义 Sobolev 空间中函数的 Fourier 变换. 对属于某个特定空间的函数, 关于其变换知道些什么?
\item 不同 $p$ 值的 Sobolev 空间之间存在哪些嵌入关系? $k$ 必须多大, 才能使 $W^{k,p}(\Omega)$ 由连续函数组成?
\item 紧性定理: 回顾 Arzela-Ascoli 定理, 它表明, 一列一致有界且导数一致有界的连续函数有一个收敛子序列. 能否将导数的一致界替换为 $L^p$ 界?
\item 迹定理: 由于 $\mathcal{D}(\Omega)$ 在 $L^p$ 中稠密, 讨论任意 $L^p$ 函数的边界值并无意义. 然而, $\mathcal{D}(\Omega)$ 一般在 $W^{k,p}(\Omega)$ 中不稠密. 对 Sobolev 空间中函数的边界值能作何断言? 如何用边界处的行为刻画 $W_0^{k,p}(\Omega)$?
\end{enumerate}
\item 然后用一节讨论 Sobolev 空间的对偶空间, 即所谓 ``负阶'' 空间.
\item 最后用一节技术性结果作结. 这些结果本身很有意义, 但将其收入本章主要是为了证明前面各节的一些主要结果. 所讨论的技术性问题如下.
\begin{enumerate}
\item 稠密性定理: $C^k(\overline\Omega)$ 中的函数在 $W^{k,p}(\Omega)$ 中稠密吗? 若 $\Omega$ 无界, 具有有界支集的函数稠密吗?
\item 坐标变换: Sobolev 空间中的函数在变换到不同区域时如何变化?
\item 延拓定理: 若 $f\in W^{k,p}(\Omega)$, 是否存在 $F\in W^{k,p}(\mathbb{R}^m)$, 使 $F|_{\Omega}=f$?
\end{enumerate}
\end{itemize}

\subsection{基本定义}\label{sec:sobolev-basic-definitions}

\begin{Definition}
\label{def:sobolev-positive-integer-spaces}
设 $k$ 为非负整数, $1\leq p\leq\infty$. 定义 $W^{k,p}(\Omega)$ 为所有满足当 $|\alpha|\leq k$ 时 $D^\alpha u\in L^p(\Omega)$ 的分布 $u\in L^p(\Omega)$ 的集合. 在 $W^{k,p}(\Omega)$ 中, 用下式定义范数:
\begin{equation}\label{eq:sobolev-positive-integer-norm-definition}
\begin{aligned}
\|u\|_{k,p}^p&:=\sum_{|\alpha|\leq k}\|D^\alpha u\|_p^p,\quad p<\infty,\\
\|u\|_{k,\infty}&:=\max_{|\alpha|\leq k}\|D^\alpha u\|_{\infty},
\end{aligned}
\end{equation}
对 $p=2$, 用下式定义内积:
\begin{equation}\label{eq:sobolev-hilbert-inner-product-definition}
(u,v)_k:=\sum_{|\alpha|\leq k}\int_{\Omega}D^\alpha\overline{u(\bm{x})}D^\alpha v(\bm{x})\,\mathrm{d}\bm{x}.
\end{equation}
对 $W^{k,2}(\Omega)$ 也使用记号 $H^k(\Omega)$.
\end{Definition}
